\documentclass[11pt,a4paper]{article}
\usepackage[margin=2.4cm]{geometry}
\usepackage{graphicx}
\usepackage{booktabs}
\usepackage{amsmath}
\usepackage[hidelinks]{hyperref}
\usepackage{microtype}
\usepackage{xcolor}

\newcommand{\wip}[1]{\textcolor{gray}{\textit{#1}}}
\title{Module E (WIP): a neutral null for ancestry sorting, and what the empirical
population structure already rules out}
\author{}
\date{Work in progress --- \today}

\begin{document}
\maketitle
\vspace{-2em}

\begin{abstract}
\noindent
This is a working note on Module E: a recombination-matched, haplodiploid
\emph{neutral null} for the 20 replicate \textit{F.\ polyctena}\,$\times$\,\textit{F.\
aquilonia} hybrid populations, built to test whether the directional ancestry
sorting seen in Modules A--C exceeds drift. The null conditions on real phased
parental haplotypes, uses the empirical recombination map, is haplodiploid, and
contains no selection. The central design choice is to stratify markers by
DiagnosticIndex (DI): low-DI markers cannot be driven by ancestry sorting and so
serve as an \emph{internal demographic control} in the same analysis that tests
the high-DI markers. Two results now stand largely on the empirical data alone,
independent of the simulation calibration. First, among-population differentiation
scales steeply with DI in the data ($F_{ST}$: $\sim$0 at neutral markers $\to$
0.31 at diagnostic markers) whereas neutral drift, being DI-blind, is flat.
Second, decomposing the diagnostic differentiation shows that a single
per-population admixture proportion (``variable founding admixture'') explains
only $\sim$6\%: the differentiation is 78\% \emph{locus-specific}. Because neutral
drift is excluded (neutral $F_{ST}\approx0$) and variable admixture is excluded
(6\%), what remains is a large, locus-specific, \emph{divergent} ancestry signal
standing against a near-zero neutral expectation. A wrinkle worth flagging: this
signal is divergent (populations differing at diagnostic loci), which is distinct
from --- though compatible with --- the parallel directional sorting of Module~A.
The definitive test, a $p$-value from the sorting statistic run through the
calibrated null, is in progress.
\end{abstract}

\section{The null model}

Rather than simulating the parental species' history, we \emph{condition} on it:
founders are drawn from real phased parental haplotypes (15 diploid females per
species, imputed and phased within species with Beagle~5.4; pool of 30
\textit{aquilonia} haplotypes and 13 \textit{polyctena} diploid females), the
mutation rate is zero, and only the post-contact hybrid phase is simulated. This
inherits true allele frequencies, true between-species divergence, and true
within-species haplotype structure. The engine is SLiM~5, non-Wright--Fisher,
haplodiploid (haploid males from unfertilised eggs, spermathecal storage, Pamilo
mating, age-structured life table), over the 26 assembled chromosomes with their
empirical cM maps (converted with \texttt{scale}$=10^{-8}$). All selection is
removed, so drift and recombination on the empirical map are the only forces. Each
parameter set runs as 20 independent demes sampled every 20 generations to
generation~160.

\paragraph{Requirements for a valid comparison.}
Three points are load-bearing and worth recording. (i) \textbf{Sample-size
matching}: composite $r^2$ and the sorting classification are both sample-size
sensitive, so simulated demes are subsampled to the empirical per-population sizes
(3--20 diploids). (ii) \textbf{Symmetric treatment}: any filter applied to the data
is applied to the null (``sorted'' \emph{is} high differentiation, so filtering the
data alone would move it toward the null mechanically). (iii) \textbf{$F_{ST}$
estimator}: with per-population $n$ of 3--20, the Nei estimator is upward-biased by
$\mathcal{O}(1/2n)$; the Hudson estimator corrects this and is used for the
structure results below. This correction matters (Section~\ref{sec:struct}).

\paragraph{Neutral alternatives excluded as controls.}
Before interpreting any excess we quantified and excluded several neutral
explanations. \emph{Founder heterogeneity} (populations founded from different
local sub-pools): even the absolute ceiling --- each deme founded from a single
\textit{aquilonia} haplotype and a single \textit{polyctena} female --- gives an
at-founding among-population $F_{ST}$ of only 0.090, because at diagnostic markers
both species are near-fixed for opposite alleles and a 50/50 admixture washes out
within-species structure. \emph{Colony (nest) sampling}: measured within-nest
excess genotype correlation is only 0.047 (polygyny), so sibling structure is a
minor effect. \emph{Mixed population ages}: age mixtures interpolate along the
LD-versus-drift trade-off and match no summary. These are recorded so they need
not be re-litigated.

\section{Differentiation and sorting scale with DI in the data, not under drift}

Markers were binned by DI over the parent-MAF$\ge$0.15 universe (40{,}000 markers,
4{,}000 per bin across ten fixed DI bins). Figure~\ref{fig:dose} and
Table~\ref{tab:dose} give the result: observed among-population $F_{ST}$ rises from
the neutral background to the most ancestry-informative markers, while the neutral
null is essentially flat. The fraction of sorted loci shows the same pattern more
sharply, from a $\sim$1.5-fold excess at low DI to $>$100-fold at high DI. Because
drift does not distinguish informative from uninformative markers, a DI-scaled
discrepancy cannot be produced by demographic misspecification --- the comparison
is internally controlled because both datasets use the same markers, MAF gate, and
frequency spectrum. (Table~\ref{tab:dose} reports Nei $F_{ST}$; the absolute
neutral value there, 0.05, is the small-sample bias corrected in
Section~\ref{sec:struct}, but the \emph{scaling} is unaffected.)

\begin{table}[htbp]\centering
\caption{Observed vs simulated statistics by DI bin at the anchor-matched neutral
cell. ``sorted'' is the fraction of loci classed as directionally sorted.}
\label{tab:dose}\small
\begin{tabular}{lrrrrrr}
\toprule
DI bin & $F_{ST}$ obs & $F_{ST}$ null & $\pi$ obs & $\pi$ null & sorted obs & sorted null\\
\midrule
$\le-90$ & 0.051 & 0.191 & 0.448 & 0.381 & 0.026 & 0.021\\
$(-75,-60]$ & 0.303 & 0.183 & 0.260 & 0.314 & 0.350 & 0.233\\
$(-40,-30]$ & 0.305 & 0.167 & 0.277 & 0.364 & 0.260 & 0.072\\
$(-25,-20]$ & 0.310 & 0.160 & 0.297 & 0.398 & 0.150 & 0.015\\
$>-15$ & 0.315 & 0.156 & 0.304 & 0.415 & 0.116 & 0.000\\
\bottomrule
\end{tabular}
\end{table}

\begin{figure}[htbp]\centering
\includegraphics[width=0.8\textwidth]{../Figures/moduleE_fig1_dose_response.pdf}
\caption{Among-population differentiation (top) and the fraction of sorted loci
(bottom, log scale) against DI bin, observed vs neutral null. Neutral drift is
DI-blind and produces an essentially flat profile; the data concentrate both
differentiation and sorting in ancestry-informative markers.}
\label{fig:dose}
\end{figure}

\section{Population structure: neutral $\approx0$, diagnostic strong}
\label{sec:struct}

This is the pivotal result. At \emph{strictly} neutral markers ($\mathrm{DI}\le-90$,
13{,}978 markers), the pairwise Hudson $F_{ST}$ among the 20 populations is
essentially \textbf{zero} (median 0.000, max 0.012), and the neutral PCA is a single
overlapping cloud with no population structure (Fig.~\ref{fig:struct}, left and
middle). At diagnostic markers ($\mathrm{DI}>-25$) the same populations are
\textbf{strongly structured}, $F_{ST}=0.31$, resolving into distinct clusters
(Fig.~\ref{fig:struct}, right).

Two consequences follow. First, the earlier ``neutral $F_{ST}\approx0.05$'' was
small-sample bias in the Nei estimator; the Hudson estimate is $\approx0$. The same
bias inflated the \emph{simulated} Nei values equally, so the sim-versus-data gap in
Table~\ref{tab:dose} survives, but the true neutral differentiation is $\approx0$.
Second, and more importantly, this dissociation cannot be drift: drift is DI-blind,
so it cannot differentiate diagnostic markers while leaving neutral markers
undifferentiated. The populations do not behave as 20 independently drifting demes
at neutral loci (independent foundings accumulate $F_{ST}\approx t/2N_e>0$); they
look effectively panmictic / very-large-$N_e$ / connected at neutral loci, while
differing strongly in ancestry. A Mantel-style test confirms the two are
dissociated: Spearman correlation between pairwise neutral $F_{ST}$ and pairwise
$|$ancestry difference$|$ is only 0.08.

\begin{figure}[htbp]\centering
\includegraphics[width=\textwidth]{../Figures/moduleE_structure.pdf}
\caption{Structure among the 20 hybrid populations. Left: pairwise neutral-marker
($\mathrm{DI}\le-90$) $F_{ST}$, Hudson estimator, clustered --- essentially zero
everywhere. Middle: individual PCA on neutral markers --- one overlapping cloud.
Right: individual PCA on diagnostic markers --- strong, structured differentiation.}
\label{fig:struct}
\end{figure}

\section{The diagnostic differentiation is locus-specific, not variable admixture}

If the populations were simply founded at different aquilonia:polyctena ratios, each
would carry a single ancestry proportion $a_j$ and its diagnostic allele
frequencies would follow $p_{ij}=f^{\mathrm{pol}}_i + a_j\,\Delta p_i$ (rank-1). We
fit this per population and decomposed the among-population variance at diagnostic
markers into: sampling (binomial), the single ancestry axis, and a genuine
locus-specific residual. Of the observed diagnostic $F_{ST}=0.313$:
\begin{center}\small
\begin{tabular}{lr}
\toprule
component & $F_{ST}$ (fraction)\\
\midrule
sampling noise & 0.049 \ (16\%)\\
single ancestry axis (variable admixture) & 0.019 \ (6\%)\\
locus-specific residual & 0.245 \ (78\%)\\
\bottomrule
\end{tabular}
\end{center}
Per-population best-fit ancestry does vary (0.35--0.65, SD 0.077;
Fig.~\ref{fig:decomp} top), but that single knob explains only $\sim$6\% of the
differentiation. Variable founding admixture is therefore rejected as the
explanation.

Combining the two structure results: drift is excluded (neutral $F_{ST}\approx0$)
and variable admixture is excluded (6\%), so the 0.245 locus-specific diagnostic
differentiation is neither. This is where the internal control pays off: a neutral
null calibrated to reproduce neutral $F_{ST}\approx0$ has essentially no drift and
therefore predicts $\approx0$ diagnostic differentiation, making the observed
0.245 \emph{almost entirely excess over neutral}.

\paragraph{An honest wrinkle: divergent, not parallel.}
The among-population \emph{variance} decomposed here measures \emph{divergence} ---
populations differing at diagnostic loci. Parallel sorting (all populations fixing
the same allele at a locus, Module~A's DI-governs-direction result) produces
agreement, i.e.\ \emph{low} among-population variance, and lives in the \emph{mean}
shift rather than this residual. So the 0.245 is a \emph{divergent} locus-specific
signal, and the residual PCA is correspondingly diffuse and high-dimensional
(Fig.~\ref{fig:decomp} bottom; PC1/PC2 $\approx$9\% each, though small per-population
$n$ confounds this view). Divergent ancestry differentiation without drift points
toward local/divergent processes (the extrinsic/climate arm, Module~C) more than a
single shared intrinsic incompatibility axis. It does not contradict Module~A ---
a parallel mean direction and divergent variance can coexist --- but it says the
\emph{differentiation among populations} is dominated by locus-specific divergence.

\begin{figure}[htbp]\centering
\includegraphics[width=0.72\textwidth]{../Figures/moduleE_ancestry_decomp.pdf}
\caption{Top: per-population best-fit aquilonia ancestry (the single admixture
axis), 0.35--0.65. Bottom: PCA of the residual after removing that axis; structure
here would be locus-specific sorting beyond variable admixture. It is diffuse and
high-dimensional (sampling-confounded at these sample sizes).}
\label{fig:decomp}
\end{figure}

\section{Supporting diagnostic: the LD landscape}

Local LD ($ld_w$, median $r^2$ in a 100~kb window) along the genome
(Fig.~\ref{fig:ldw}) shows that the null reproduces the recombination-driven LD
\emph{landscape} well --- peaks co-localise on the same chromosomes, per-chromosome
Spearman 0.94--0.95, robust across generations --- but carries a diffuse genome-wide
LD \emph{floor} the data lack (mean $\sim$0.17--0.21 vs 0.09), decaying only slowly
from gen20 to gen160. A within-group analysis (removing the between-population
component) showed the founders themselves carry within-species LD matching the
empirical within-population LD, so the excess floor is not founder-inherited; it is
the between-population (Wahlund) component, i.e.\ the same over-differentiation
(too-small $N_e$) seen in the $F_{ST}$ analysis. This independently corroborates
that the null over-differentiates the demes and motivates the large-$N_e$
calibration.

\begin{figure}[htbp]\centering
\includegraphics[width=\textwidth]{../Figures/moduleE_ldw_manhattan.pdf}
\caption{Local LD along the genome, empirical vs neutral null over generations
(K=6250, pooled, matched $N$). The recombination-driven landscape is reproduced
(peaks co-localise); the null carries a diffuse excess floor between peaks --- the
between-population (Wahlund) component of the over-differentiation.}
\label{fig:ldw}
\end{figure}

\section{Where this stands, and next steps \wip{(WIP)}}

Two demographic explanations for the low-neutral / high-diagnostic dissociation are
now excluded from the empirical data alone: \emph{drift} (neutral $F_{ST}\approx0$,
and drift is DI-blind) and \emph{variable founding admixture} (6\% of the diagnostic
differentiation). What remains is a large ($F_{ST}\approx0.245$), locus-specific,
\emph{divergent} ancestry signal against a near-zero neutral expectation. The
simulation's role is now sharply defined: calibrate the null to reproduce neutral
$F_{ST}\approx0$ (the very-large-$N_e$ / effectively-panmictic-at-neutral regime),
at which point the diagnostic excess is read against $\approx0$.

\begin{enumerate}\itemsep2pt
\item \textbf{Calibrated-null $p$-value.} Run the sorting statistic
(\texttt{parallelism\_stats}) through the null calibrated on the neutral markers,
extending carrying capacity upward until neutral $F_{ST}\approx0$. This puts a
$p$-value on the excess and separates the parallel (mean) from the divergent
(variance) components properly. \wip{Higher-$K$ simulations are running on a
larger machine.}
\item \textbf{Divergent vs parallel.} Characterise the residual structure with
individual-level data (controlling sampling), and test whether it aligns with
climate (reconciling with Modules~A and~C: parallel direction vs divergent
differentiation).
\item \textbf{Non-independence.} Neutral $F_{ST}\approx0$ is formally degenerate
between enormous $N_e$ and gene flow / shared ancestry among ``populations'';
extra within-event sampling and demographic modelling (cf.\ Nouhaud et al.) would
separate these, and it bears on how many independent replicates the sorting test
really has.
\end{enumerate}

\noindent\wip{All analyses run from a 19.5~MB compact empirical bundle (replacing
the 1~GB dataset); scripts and the portable simulation pipeline are in
\texttt{moduleE\_slim/} in the repository.}

\end{document}
